\reportsection{Hiding}
\label{sec:hiding}

A zero-knowledge protocol relies on more than hiding, and the book's stronger requirement is
\emph{equivocation}, which a Merkle commitment satisfies in the random-oracle model through
an inefficient simulator that programs the oracle, though never in the \emph{trapdoor} form
an efficient simulator would need, since there is no $\Setup$ secret to program
($\td=\bot$)~\cite{CY26}.  We work with the indistinguishability form of hiding below, which
rests entirely on the salt space, and show that the current salt is too small to carry it.

\begin{definition}[Hiding]\label{def:vc-hiding}
A randomised $\VC$ whose $\Commit$ draws salts from a space $\rndspace$ is
\emph{hiding} if, for any two messages of length $\MTMessageLength$ and any
adversary that after seeing $\comm$ adaptively opens positions on which the
messages agree, the advantage in deciding which was committed is at most
$\negl$.  For Merkle this refines to \emph{root hiding}, meaning the real root is
statistically close to one simulated without $\msgvec$, strengthened to
\emph{selective-opening hiding}, which must also hide through the disclosed
salts and authentication paths of partial openings.
\end{definition}

In a Merkle commitment scheme, we refer to salt simply as the random per-leaf nonce.
Each leaf is itself a basic commitment $H(\msgvec_i,\text{salt})$, and the salt travels
in the authentication path beside the sibling digests so the verifier can recompute
the leaf.  Opening one leaf commitment under two different salts would force a
hash collision, so the salt binds a value to its position even as it hides it.
Hiding then rests entirely on the salt space, which must be both well sampled
and large enough to resist search within the security budget, neither of which a
deterministic salt satisfies, since the receiver can recompute candidate leaf commitments.

The digest and the salt scale differently in the query budget, which is what
\cref{eq:rule} is tracking when it derives $D$ from $h$ and derives $k$ and $S$ from
$\secpar$ alone.  A binding adversary wins on any collision among its own queries, so its
advantage grows as $\querybudget^2/2^{\digestbits}$ and the digest has to carry $2h$ bits of
headroom.  A hiding adversary has to hit one particular salted leaf input, so its advantage
grows as $\msglen\querybudget/2^{\saltbits}$, and the salt needs only $\secpar$ bits.

\begin{proposition}[Salt cardinality]\label{prop:salt-card}
Let the salt be drawn from $k$ elements of a field with $b=\lfloor\log_2|\field|\rfloor$
bits per element, or from $S$ bytes.  Then $|\rndspace|\ge2^{\secpar}$ whenever
$k\ge\lceil\secpar/b\rceil$ or $S\ge\lceil\secpar/8\rceil$.
\status{Lean-backed}
\end{proposition}

This is what \cref{eq:rule} computes for $k$ and $S$ and the one part of hiding that is
machine-checked, while the bound it feeds into remains an \emph{intended reduction}
(\cref{app:hiding-reduction}).

Looking again at the Arkworks library, the current hiding hasher repeats the
undersized-parameter mistake we saw for binding, this time in the salt space
(\cref{app:hasher-trait}): the salt is one field element for every field and every target
$\lambda$, drawn uniformly, so the sampling requirement is met and what falls short is the
number of values it can take.

Take a target security parameter $\lambda=128$ and BabyBear again.  Each field element
carries at most $b=\lfloor\log_2 p\rfloor=30$ bits under the conservative accounting of
\cref{eq:rule}, so a one-element salt ranges over roughly $2^{30}$ values.  An
adversary who knows the two candidate messages recovers which was committed by hashing
$2^{30}$ candidate leaves against the opened position.  Reaching $2^{\lambda}$ instead needs
$k=\lceil\lambda/b\rceil=\lceil128/30\rceil=5$ elements.  The same constructor gives
$k=3$ for Goldilocks at $b=63$, and $k=1$ for BN254 at $b=253$, where a single element
already carries more than the target.
